\section{Estructura del proyecto}

El proyecto se organiza en tres directorios principales equivalentes a los tres entornos descritos en la sección anterior. Esto asegura el aislamiento de dependencias y la claridad en la lectura dentro del repositorio.

\begin{forest}
	for tree={
		font=\ttfamily,
		grow'=0,
		folder,           
		l sep=1.5em,  
		s sep=2pt,        
	}
	[Simulator/
	[\fitem{core/}{\# C++17/CUDA + bindings pybind11}]
	[\fitem{GUI/}{\# Aplicación PyQt5}]
	[\fitem{training/}{\# Scripts de entrenamiento en PyTorch}]
	]
\end{forest}

\subsection{Directorio core}

Contiene la implementación de bajo nivel del autómata celular y del autómata celular neuronal.

\begin{table}[H]
    \centering
    \begin{tabular}{c|p{0.52\linewidth}}
        \hline
        \textbf{Archivo} & \textbf{Responsabilidad} \\ \hline
        \texttt{ca.h} / \texttt{ca.cu} & Definición e implementación del autómata celular clásico en CUDA. \\ \hline
        \texttt{nca.h} / \texttt{nca.cu} & Definición e implementación del autómata celular neuronal. \\ \hline
        \texttt{bindings/python\_module.cpp} & Capa pybind11 que expone las clases \texttt{CA} y \texttt{NCA} al intérprete de Python. \\ \hline
        \texttt{CMakeLists.txt} & Configuración del \textit{build} con CMake. \\ \hline
        \texttt{bind.sh} & \textit{Script} de automatización. \\ \hline
        \texttt{requirements.txt} & Dependencia única: \texttt{pybind11}. \\ \hline
        \texttt{build/} & Directorio generado por CMake (no versionado). \\ \hline
        \texttt{venv\_bind/} & Entorno virtual local para el proceso de compilación. \\ \hline
    \end{tabular}
    \caption{Contenido del directorio \texttt{core/}.}
\end{table}

\subsection{Directorio GUI}

Implementa la aplicación gráfica completa siguiendo el patrón Modelo-Vista-Controlador descrito en la sección \ref{sec:arquitectura}.

\begin{table}[H]
    \centering
    \begin{tabular}{ p{0.4\textwidth} | p{0.54\textwidth}}
        \hline
        \textbf{Ruta} & \textbf{Responsabilidad} \\ \hline
        \texttt{main.py} & Instancia \texttt{QApplication} y \texttt{MainWindow}. \\ \hline
        \texttt{models/}  \texttt{cellularAutomaton.py} & \textit{Wrapper} de Python sobre la clase \texttt{CA} expuesta por pybind11. \\ \hline
        \texttt{models/}  \texttt{neuralCellularAutomaton.py} & \textit{Wrapper} de Python sobre la clase \texttt{NCA}. \\ \hline
        \texttt{controllers/}  \texttt{SimulationController.py} & Controlador principal. Coordina los \textit{widgets}, el modelo y el \textit{prefetch} asíncrono. \\ \hline
        \texttt{controllers/}  \texttt{SimulationConfig.py} & Encapsulamiento de la configuración del autómata. \\ \hline
        \texttt{ui/}  \texttt{MainWindow.py} & Ventana principal. Contiene un \texttt{QTabWidget} con las pestañas de simulación y entrenamiento. \\ \hline
        \texttt{ui/}  \texttt{ControlPane.py} & Panel izquierdo de la pestaña de simulación: reproducción, paso, retroceso, borrado, aleatorización y velocidad. \\ \hline
        \texttt{ui/}  \texttt{SimulationWidget.py} & \textit{Widget} central. Renderiza la cuadrícula con \texttt{QImage} en formato indexado y emite señales de \textit{click}. \\ \hline
        \texttt{ui/}  \texttt{ConfigPane.py} & Panel derecho. Configura vecindad, estados, frontera y regla de transición. \\ \hline
        \texttt{ui/}  \texttt{NeighbourWidget.py} & Editor visual de la vecindad. Soporta múltiples grupos. \\ \hline
        \texttt{ui/}  \texttt{TrainingPane.py} & Pestaña de entrenamiento. Captura hiperparámetros, lanza \texttt{QProcess} y muestra el \textit{log} en tiempo real. \\ \hline
    \end{tabular}
    \caption{Contenido del directorio \texttt{GUI/}.}
\end{table}

\subsection{Directorio training}

Reúne los \textit{scripts} de entrenamiento del autómata celular neuronal, las evoluciones de referencia utilizadas como conjuntos de datos y los modelos entrenados a partir de ellas.

\begin{table}[H]
    \centering
    \begin{tabular}{p{0.2\linewidth}|p{0.7\linewidth}}
        \hline
        \textbf{Archivo /}   \quad \textbf{directorio} & \textbf{Responsabilidad} \\ \hline
        \texttt{scripts/} \texttt{main.py} & \textit{Script} principal. Implementa la extracción de pares (vecindad, objetivo), la construcción dinámica del MLP en PyTorch y el ciclo de entrenamiento. \\ \hline
        \texttt{scripts/} \texttt{test.py} & Implementa en NumPy el paso del autómata celular neuronal y compara su evolución con la regla exacta del juego de la vida. \\ \hline
        \texttt{scripts/} \texttt{png\_to\_\allowbreak evolution.py} & Convierte una secuencia de imágenes PNG segmentadas en una evolución con el formato \texttt{(T, H, W)} en \texttt{.npy}/\texttt{.mat}. \\ \hline
        \texttt{CA/} & Evoluciones de referencia, una carpeta por experimento: las reglas del juego de la vida (\texttt{evo} para entrenamiento, \texttt{base} como condición de prueba), \texttt{Sandpile/}, \texttt{LangtonAnt/}, \texttt{FLAME3/}, \texttt{SEVIR/} y \texttt{Physarum/}. \\ \hline
        \texttt{NCA/} & Pesos entrenados (\texttt{.npz}) por experimento, las evoluciones que producen (\texttt{t\_*}) y los resultados de los barridos; \texttt{ga\_runs/} contiene la búsqueda genética. \\ \hline
        \texttt{notebooks/} & Un \textit{notebook} por experimento: de Python (\texttt{.ipynb}) para convertir los datos externos en evoluciones (FLAME~3, SEVIR y \textit{Physarum}) y de Wolfram Mathematica (\texttt{.nb}) para el análisis y las figuras (juego de la vida, \textit{sandpile}, hormiga de Langton y FLAME~3). \\ \hline
        \texttt{Dataset/} & Datos crudos descargados (no versionado). \\ \hline
    \end{tabular}
    \caption{Contenido del directorio \texttt{training/}.}
\end{table}

\subsection{Flujo de artefactos entre directorios}

El flujo de artefactos entre directorios es el siguiente:

\begin{enumerate}
    \item El \textit{build} en \texttt{core/build/} produce \texttt{CA.cpython-*.so} (con las clases \texttt{CA} y \texttt{NCA}), que se copia al directorio \texttt{GUI/}.
    \item La GUI registra evoluciones mediante el botón de grabación de \texttt{ControlPane} y las exporta a \texttt{.npy}/\texttt{.mat} en cualquier ubicación, típicamente \texttt{training/CA/}.
    \item El \textit{script} de entrenamiento lee dichos archivos y produce un \texttt{.npz} con pesos, sesgos y metadatos.
    \item La GUI carga el \texttt{.npz} desde \texttt{SimulationController} y cambia al modo NCA para reproducir la evolución aprendida.
\end{enumerate}

Este flujo unidireccional garantiza que los tres entornos puedan evolucionar de forma independiente: un cambio en el módulo nativo no requiere reentrenar, y un cambio en el entrenamiento no requiere recompilar.
